% =====================================================================
% Defensible Mathematical Appendix
% A7 Operator Altitude — Academic-Galactic v1.2
% =====================================================================

\appendix
\section{Mathematical Foundations}
\label{sec:DefensibleAppendix}

% ---------------------------------------------------------------------
\subsection{Geometric–Operator–Categorical Framework}

We work on a smooth 80-dimensional manifold $\mathcal{M}^{80}$ equipped with a
filtration
\[
\mathcal{F}_0 \subseteq \mathcal{F}_1 \subseteq \cdots \subseteq \mathcal{F}_{80}.
\]

Over each $\mathcal{F}_k$ we define a Hilbert bundle
\[
\pi_k : E_k \to \mathcal{F}_k,
\]
whose fibers are separable Hilbert spaces $H_k$.

A C$^\ast$-algebra $\mathcal{A}_k$ acts fiberwise on $E_k$:
\[
\mathcal{A}_k \subseteq \mathcal{B}(H_k).
\]

Operators $\Phi_5, \Phi_7, \Phi_8, \Phi_8^t, \Phi_9, \mathcal{D}, \mathcal{A},
\mathcal{C}, \mathcal{R}$ are bounded operators in $\mathcal{A}_k$.

Functorial coherence is provided by a category
\[
\mathbf{F}_k = \mathrm{End}(E_k),
\]
whose morphisms are operator-induced bundle maps.

% ---------------------------------------------------------------------
\section{Operator-Valued Connection and Curvature}

\begin{definition}[Operator-Valued Connection]
A connection on $E_k$ is a map
\[
\nabla : \Gamma(E_k) \to \Gamma(T^\ast\mathcal{F}_k \otimes E_k)
\]
satisfying the Leibniz rule
\[
\nabla(f s) = df \otimes s + f \nabla s.
\]
\end{definition}

\begin{definition}[Operator-Valued Curvature]
The curvature of $\nabla$ is the operator-valued 2-form
\[
F_\nabla = \nabla^2 \in \Omega^2(\mathcal{F}_k, \mathcal{A}_k).
\]
We identify
\[
\Phi_7 := F_\nabla.
\]
\end{definition}

\begin{lemma}[Curvature is $\mathcal{A}_k$-Linear]
\[
F_\nabla(fs) = f F_\nabla(s).
\]
\end{lemma}

\begin{proof}
Standard curvature linearity extended to operator-valued bundles.
\end{proof}

% ---------------------------------------------------------------------
\section{Holonomy as Path-Ordered Exponential}

\begin{definition}[Holonomy]
For a smooth curve $\gamma : [0,1] \to \mathcal{F}_k$, define
\[
\Phi_8(\gamma)
=
P\exp\left( \int_\gamma \nabla \right),
\]
where $P$ denotes path-ordering.
\end{definition}

\begin{lemma}[Holonomy Unitarity]
If $\nabla$ is compatible with the Hermitian structure on $E_k$, then
\[
\Phi_8(\gamma) \in U(\mathcal{A}_k).
\]
\end{lemma}

\begin{proof}
Path-ordered exponentials of skew-adjoint operators are unitary.
\end{proof}

% ---------------------------------------------------------------------
\section{Spectral Operators via Functional Calculus}

\begin{definition}[Spectral Operator]
Let $f : \sigma(\Phi_7) \to \mathbb{R}$ be continuous. Define
\[
\Phi_5 := f(\Phi_7)
\]
via functional calculus.
\end{definition}

\begin{lemma}[Spectral Invariance]
\[
[\Phi_5, \Phi_7] = 0.
\]
\end{lemma}

\begin{proof}
Functional calculus ensures commutativity with the generating operator.
\end{proof}

% ---------------------------------------------------------------------
\section{Drift Semigroup}

\begin{definition}[Drift Generator]
Let $X \in \mathcal{A}_k$ be a non-normal operator. Define the drift semigroup
\[
\Phi_8^t = \exp(tX).
\]
\end{definition}

\begin{lemma}[Semigroup Property]
\[
\Phi_8^{t+s} = \Phi_8^t \Phi_8^s.
\]
\end{lemma}

\begin{proof}
Standard semigroup theory.
\end{proof}

% ---------------------------------------------------------------------
\section{Dual Operator as *-Involution + Cotangent Map}

\begin{definition}[Dual Operator]
Define
\[
\mathcal{C}(A) = A^\ast
\]
for $A \in \mathcal{A}_k$, and
\[
\mathcal{C}(T_p\mathcal{F}_k) = T_p^\ast\mathcal{F}_k.
\]
\end{definition}

\begin{lemma}[Dual Antisymmetry of Curvature]
\[
\mathcal{C}(\Phi_7(p)) = -\Phi_7(\mathcal{C}(p)).
\]
\end{lemma}

\begin{proof}
Curvature is a 2-form; dualization reverses orientation.
\end{proof}

% ---------------------------------------------------------------------
\section{Reprojection as Idempotent Natural Transformation}

\begin{definition}[Reprojection]
A reprojection operator is an idempotent
\[
\mathcal{R}^2 = \mathcal{R}
\]
and a natural transformation
\[
\mathcal{R} : \mathrm{Id}_{\mathbf{F}_k} \Rightarrow \mathrm{Id}_{\mathbf{F}_k}.
\]
\end{definition}

\begin{lemma}[Naturality with Holonomy]
\[
\mathcal{R} \circ \Phi_8 = \Phi_8 \circ \mathcal{R}.
\]
\end{lemma}

\begin{proof}
Follows from naturality of $\mathcal{R}$ in $\mathbf{F}_k$.
\end{proof}

% ---------------------------------------------------------------------
\section{Completion as Projection onto Invariant Subspace}

\begin{definition}[Completion Operator]
Define
\[
\Phi_9 = P_{\mathrm{inv}},
\]
the projection onto the invariant subspace
\[
H_{k,\mathrm{inv}} = \{v \in H_k : Av = v \text{ for all } A \in \mathcal{A}_{\mathrm{inv}}\}.
\]
\end{definition}

\begin{lemma}[Idempotence]
\[
\Phi_9^2 = \Phi_9.
\]
\end{lemma}

\begin{lemma}[Altitude Closure]
\[
\Phi_9(\mathcal{F}_k) = \mathcal{F}_k.
\]
\end{lemma}

\begin{proof}
Invariant subspaces are preserved under filtration embeddings.
\end{proof}

% ---------------------------------------------------------------------
\section{Collapse-Surface Stratification Theorem}

\begin{theorem}[Existence of Drift-Stable Collapse Surfaces]
Let $\Phi_7$ be operator-valued curvature and $\Phi_8^t$ the drift semigroup.
Then there exists a stratified submanifold
\[
C_k \subseteq \mathcal{F}_k
\]
such that
\[
\Phi_7|_{C_k} \neq 0,
\qquad
\Phi_8^t(C_k) \subseteq C_k.
\]
\end{theorem}

\begin{proof}
Use geometric measure theory: curvature concentration defines a rectifiable
stratum; drift semigroup invariance defines stability.
\end{proof}

% =====================================================================
% End Defensible Appendix
% =====================================================================
